\chapter{Kommentierte Quellen}

Achtundzwanzig Werke, eingeordnet: Zu jedem steht, wofür man es aufschlägt und was es leistet, das andere nicht leisten.

\lit{1}{Aumann, R., Hart, S. (Hg.): \textit{Handbook of Game Theory with Economic Applications}. Band I--IV, Elsevier 1992--2015.}{Das Nachschlagewerk des Fachs. Jedes Kapitel ist eine Übersichtsarbeit einer führenden Person zu ihrem Gebiet. Kein Lehrbuch, sondern die Stelle, an der man nachsieht, was zu einer Frage bekannt ist und wer es gezeigt hat.}

\lit{2}{Axelrod, R.: \textit{The Evolution of Cooperation}. Basic Books 1984.}{Das Buch zu Kapitel~8, geschrieben für ein breites Publikum und dennoch sorgfältig. Enthält beide Turniere, die vier Eigenschaften erfolgreicher Strategien und Anwendungen vom Stellungskrieg bis zur Bakterienökologie. \textbf{Die spätere Kritik --- Rauschen, Turniermanipulation --- muss man dazulesen.}}

\lit{3}{Berninghaus, S., Ehrhart, K.-M., Güth, W.: \textit{Strategische Spiele}. 3.\ Aufl., Springer 2010.}{\textbf{Auf Deutsch und mit Blick auf Experimente.} Güth ist der Erfinder des Ultimatumspiels, und das Buch verbindet die Theorie durchgehend mit dem, was in Versuchen tatsächlich geschieht. Die passende Ergänzung zu Kapitel~15.}

\lit{4}{Binmore, K.: \textit{Playing for Real}. Oxford 2007.}{\textbf{Das Buch zum Verstehen.} Binmore erklärt nicht nur die Sätze, sondern auch, was sie bedeuten und wo sie missbraucht werden. Eigenwillig geschrieben, streitlustig und außerordentlich klar in den Begriffen. Wer diesen Band gelesen hat und weiterwill, liest hier weiter.}

\lit{5}{Camerer, C.: \textit{Behavioral Game Theory}. Princeton 2003.}{\textbf{Die maßgebliche Darstellung dessen, was Menschen tatsächlich tun.} Enthält die Versuchsergebnisse zu allen Spielen dieses Bandes, die Modelle mit begrenzter Denktiefe und die Lerntheorien. Die Grundlage von Kapitel~15.}

\lit{6}{Dixit, A., Nalebuff, B.: \textit{The Art of Strategy}. Norton 2008.}{Die zugänglichste Einführung: ohne Formeln, mit vielen Beispielen aus Politik, Sport und Wirtschaft. Als Ergänzung geeignet; als alleinige Quelle nicht, weil die Vorbehalte zu kurz kommen.}

\lit{7}{Fudenberg, D., Tirole, J.: \textit{Game Theory}. MIT Press 1991.}{Das Standardwerk für Fortgeschrittene, mit dem Schwerpunkt auf dynamischen Spielen und unvollständiger Information. Enthält die Verfeinerungen des Gleichgewichtsbegriffs vollständig --- also das, was Kapitel~7 nur andeuten konnte.}

\lit{8}{Gibbons, R.: \textit{A Primer in Game Theory}. Harvester 1992.}{Die kürzeste brauchbare Einführung für Wirtschaftswissenschaftler --- rund zweihundert Seiten, vier Kapitel, viele durchgerechnete Anwendungen. Gut geeignet, um die Anwendungen kennenzulernen, die dieser Band aus Platzgründen ausgelassen hat.}

\lit{9}{Gintis, H.: \textit{Game Theory Evolving}. 2.\ Aufl., Princeton 2009.}{Ein Lehrbuch, das fast vollständig aus Aufgaben besteht --- über zweihundert, mit Lösungen. Verbindet klassische und evolutionäre Spieltheorie und ist die beste Quelle, wenn man rechnen üben will statt zu lesen.}

\lit{10}{Hofbauer, J., Sigmund, K.: \textit{Evolutionary Games and Population Dynamics}. Cambridge 1998.}{\textbf{Die mathematische Darstellung zu Teil~IV.} Enthält die Replikatordynamik vollständig, den Zusammenhang zu Lotka-Volterra, die Klassifikation der Bahnen und die Sätze über Erhaltungsgrößen. Anspruchsvoll und unverzichtbar, wenn man Kapitel~10 vertiefen will.}

\lit{11}{Holler, M., Illing, G., Napel, S.: \textit{Einführung in die Spieltheorie}. 8.\ Aufl., Springer Gabler 2019.}{\textbf{Auf Deutsch das Standardlehrbuch.} Deckt alle Teile dieses Bandes ab, mit Schwerpunkt auf wirtschaftswissenschaftlichen Anwendungen und einem ausführlichen Teil zu kooperativen Spielen --- der Vertiefung zu Kapitel~14.}

\lit{12}{Kagel, J., Roth, A. (Hg.): \textit{The Handbook of Experimental Economics}. Band I und II, Princeton 1995/2016.}{Die systematische Sammlung dessen, was Versuche ergeben haben --- nach Spielklassen geordnet, mit den Streitpunkten der Auswertung. Wer wissen will, wie belastbar ein bestimmter Befund ist, sieht hier nach.}

\lit{13}{Krishna, V.: \textit{Auction Theory}. 2.\ Aufl., Academic Press 2009.}{Die Vertiefung zu Kapitel~13 --- einschließlich des Beweises für das Erstpreis-Gleichgewicht, des Erlösäquivalenzsatzes in voller Allgemeinheit und der Theorie der Mehrgegenstandsauktionen.}

\lit{14}{Luce, R., Raiffa, H.: \textit{Games and Decisions}. Wiley 1957.}{Das Buch, das die Spieltheorie über ihren Gründerkreis hinaus bekannt machte. Bemerkenswert sind die begrifflichen Erörterungen: Luce und Raiffa benennen die Schwierigkeiten des Gleichgewichtsbegriffs mit einer Schärfe, die in späteren Lehrbüchern selten wurde.}

\lit{15}{Maynard Smith, J.: \textit{Evolution and the Theory of Games}. Cambridge 1982.}{Die Darstellung der evolutionären Spieltheorie durch ihren Begründer. Enthält Falke-Taube, den berichtigten Vergelter, asymmetrische Spiele und die Anwendungen auf Geschlechterverhältnis und Paarungsverhalten. Klar geschrieben und ohne überflüssige Technik.}

\lit{16}{Myerson, R.: \textit{Game Theory: Analysis of Conflict}. Harvard 1991.}{Das gründlichste Lehrbuch zum Mechanismenentwurf, von einem der Nobelpreisträger dieses Gebiets. Enthält Gibbard-Satterthwaite, das Offenbarungsprinzip und die Bayesschen Gleichgewichte --- die Fortsetzung der Kapitel~12 und 13.}

\lit{17}{Nash, J.: \textit{Equilibrium Points in N-Person Games} (1950) und \textit{Non-Cooperative Games} (1951).}{Zwei Seiten und neun Seiten. Die erste Arbeit enthält Definition und Existenz über Kakutani, die zweite den Beweis über Brouwer, den Abschnitt 6.3 wiedergibt. Beide sind kurz genug, um sie an einem Abend zu lesen.}

\lit{18}{von Neumann, J., Morgenstern, O.: \textit{Theory of Games and Economic Behavior}. Princeton 1944.}{Das Gründungswerk, sechshundert Seiten. Enthält den Minimax-Satz, die Nutzentheorie aus Satz 4.1 und die kooperative Theorie. \textbf{Man liest es heute nicht mehr durch}, aber die Einleitung über das Verhältnis von Mathematik und Wirtschaftswissenschaft lohnt.}

\lit{19}{Osborne, M., Rubinstein, A.: \textit{A Course in Game Theory}. MIT Press 1994.}{\textbf{Für die mathematische Seite die maßgebliche Quelle.} Enthält die Beweise, die dieser Band ausgelassen hat --- Brouwer, Arrow, die Unabhängigkeit der Dominanzelimination. Knapp, streng und frei verfügbar. Rubinsteins Vorbehalte gegen die Anwendung stehen im Vorwort.}

\lit{20}{Ostrom, E.: \textit{Governing the Commons}. Cambridge 1990.}{\textbf{Die empirische Widerlegung einer theoretischen Vorhersage.} Hunderte von Fällen, in denen Selbstverwaltung gemeinsamer Ressourcen über Jahrhunderte funktioniert hat, und die acht Regeln, die sie gemeinsam haben. Zu Kapitel~11 unverzichtbar.}

\lit{21}{Peters, H.: \textit{Game Theory: A Multi-Leveled Approach}. 2.\ Aufl., Springer 2015.}{Ein Lehrbuch mit ungewöhnlichem Aufbau: Es behandelt jedes Thema zweimal --- einmal elementar, einmal streng. Der kooperative Teil ist besonders gut, einschließlich des Satzes von Shapley über konvexe Spiele.}

\lit{22}{Rieck, C.: \textit{Spieltheorie --- Eine Einführung}. 14.\ Aufl., Rieck 2013.}{\textbf{Auf Deutsch, für den Einstieg.} Anschaulich, mit vielen Beispielen und ohne mathematische Vorkenntnisse lesbar. Wer den Anfang dieses Bandes zu schnell fand, findet hier den sanfteren Weg.}

\lit{23}{Roth, A.: \textit{Who Gets What --- and Why}. Houghton Mifflin 2015.}{\textbf{Die Theorie in der Anwendung, vom Verantwortlichen selbst erzählt.} Schulplätze, Assistenzarztstellen, Nierentauschketten --- einschließlich der Schwierigkeiten, die in den Aufsätzen nicht stehen. Zu Kapitel~13 die beste Ergänzung.}

\lit{24}{Rubinstein, A.: \textit{Perfect Equilibrium in a Bargaining Model}. Econometrica 50, 1982.}{Zwölf Seiten, die eine Lücke von dreißig Jahren schlossen: Warum sollte gerade die Nash-Verhandlungslösung herauskommen? Ein Musterbeispiel für die nichtkooperative Begründung eines kooperativen Begriffs.}

\lit{25}{Schelling, T.: \textit{The Strategy of Conflict}. Harvard 1960.}{\textbf{Das Buch mit den wenigsten Formeln und den meisten Ideen.} Glaubwürdige Festlegung, Selbstbindung, Brennpunkte, die Kunst des Drohens. Schelling arbeitete ohne Matrizen und hat dem Fach mehr Begriffe geliefert als die meisten Theoretiker.}

\lit{26}{Sigmund, K.: \textit{The Calculus of Selfishness}. Princeton 2010.}{Eine knappe, elegante Darstellung der Evolution von Zusammenarbeit --- direkte und indirekte Gegenseitigkeit, Bestrafung, Ruf. Verbindet die Teile III und IV dieses Bandes und ist mathematisch sauber, ohne schwerfällig zu werden.}

\lit{27}{Vickrey, W.: \textit{Counterspeculation, Auctions, and Competitive Sealed Tenders}. Journal of Finance 16, 1961.}{Die Arbeit, die den Mechanismenentwurf vorwegnahm --- zwanzig Jahre bevor das Gebiet entstand. Enthält Satz 13.1 und den Gedanken der Erlösgleichheit in einem Sonderfall.}

\lit{28}{Wiese, H.: \textit{Kooperative Spieltheorie}. Oldenbourg 2005.}{\textbf{Auf Deutsch und ausschließlich zum kooperativen Zweig.} Kern, Shapley-Wert, Verhandlungslösungen, Stimmgewichte --- mit vollständigen Beweisen und vielen Rechenbeispielen. Die Vertiefung zu Kapitel~14.}

\begin{merksatz}
\boxtitel{Wenn Sie nur drei Bücher kaufen}
\begin{center}
\renewcommand{\arraystretch}{1.6}
\begin{tabular}{@{}P{1.0cm}P{3.2cm}P{4.0cm}@{}}
\toprule
\textbf{[19]} & \textbf{Osborne, Rubinstein} & \textbf{Für die Beweise}\\
\textbf{[4]} & \textbf{Binmore} & Für das Verständnis\\
\textbf{[5]} & \textbf{Camerer} & \textbf{Für die Wirklichkeit}\\
\bottomrule
\end{tabular}
\end{center}

\begin{center}
\begin{minipage}{0.88\textwidth}
\kursiv{Osborne und Rubinstein enthalten alles, was dieser Band an Beweisen ausgelassen hat, und sind frei verfügbar. Binmore erklärt, was die Sätze bedeuten --- und was sie nicht bedeuten. Camerer steht nicht zufällig an dritter Stelle: \textbf{Ohne die Versuchsergebnisse bleibt die Theorie eine Aussage über Wesen, die es nicht gibt.}}
\end{minipage}
\end{center}

\begin{center}
\begin{minipage}{0.88\textwidth}
\kursiv{Vier Ergänzungen nach Interessenlage: [25] Schelling für die Begriffe, die keine Formel braucht; [23] Roth für die gelungene Anwendung; [20] Ostrom für die misslungene Vorhersage; [10] Hofbauer und Sigmund für Teil~IV. Auf Deutsch: [22] Rieck als Einstieg, [11] Holler und andere als Lehrbuch, [28] Wiese für den kooperativen Zweig.}
\end{minipage}
\end{center}
\end{merksatz}

\bigskip
\begin{flushright}
\itshape\color{grau}
Ein Optimum ist die beste Wahl.\\
Ein Gleichgewicht ist eine Lage,\\
aus der niemand allein ausbricht.\\[10pt]
\normalfont\color{grau}\upshape
--- der Unterschied, um den\\
\itshape es in diesem Band geht.
\end{flushright}
